% Options for packages loaded elsewhere
\PassOptionsToPackage{unicode}{hyperref}
\PassOptionsToPackage{hyphens}{url}
\documentclass[
]{article}
\usepackage{xcolor}
\usepackage[margin=1in]{geometry}
\usepackage{amsmath,amssymb}
\setcounter{secnumdepth}{-\maxdimen} % remove section numbering
\usepackage{iftex}
\ifPDFTeX
  \usepackage[T1]{fontenc}
  \usepackage[utf8]{inputenc}
  \usepackage{textcomp} % provide euro and other symbols
\else % if luatex or xetex
  \usepackage{unicode-math} % this also loads fontspec
  \defaultfontfeatures{Scale=MatchLowercase}
  \defaultfontfeatures[\rmfamily]{Ligatures=TeX,Scale=1}
\fi
\usepackage{lmodern}
\ifPDFTeX\else
  % xetex/luatex font selection
\fi
% Use upquote if available, for straight quotes in verbatim environments
\IfFileExists{upquote.sty}{\usepackage{upquote}}{}
\IfFileExists{microtype.sty}{% use microtype if available
  \usepackage[]{microtype}
  \UseMicrotypeSet[protrusion]{basicmath} % disable protrusion for tt fonts
}{}
\makeatletter
\@ifundefined{KOMAClassName}{% if non-KOMA class
  \IfFileExists{parskip.sty}{%
    \usepackage{parskip}
  }{% else
    \setlength{\parindent}{0pt}
    \setlength{\parskip}{6pt plus 2pt minus 1pt}}
}{% if KOMA class
  \KOMAoptions{parskip=half}}
\makeatother
\usepackage{color}
\usepackage{fancyvrb}
\newcommand{\VerbBar}{|}
\newcommand{\VERB}{\Verb[commandchars=\\\{\}]}
\DefineVerbatimEnvironment{Highlighting}{Verbatim}{commandchars=\\\{\}}
% Add ',fontsize=\small' for more characters per line
\usepackage{framed}
\definecolor{shadecolor}{RGB}{248,248,248}
\newenvironment{Shaded}{\begin{snugshade}}{\end{snugshade}}
\newcommand{\AlertTok}[1]{\textcolor[rgb]{0.94,0.16,0.16}{#1}}
\newcommand{\AnnotationTok}[1]{\textcolor[rgb]{0.56,0.35,0.01}{\textbf{\textit{#1}}}}
\newcommand{\AttributeTok}[1]{\textcolor[rgb]{0.13,0.29,0.53}{#1}}
\newcommand{\BaseNTok}[1]{\textcolor[rgb]{0.00,0.00,0.81}{#1}}
\newcommand{\BuiltInTok}[1]{#1}
\newcommand{\CharTok}[1]{\textcolor[rgb]{0.31,0.60,0.02}{#1}}
\newcommand{\CommentTok}[1]{\textcolor[rgb]{0.56,0.35,0.01}{\textit{#1}}}
\newcommand{\CommentVarTok}[1]{\textcolor[rgb]{0.56,0.35,0.01}{\textbf{\textit{#1}}}}
\newcommand{\ConstantTok}[1]{\textcolor[rgb]{0.56,0.35,0.01}{#1}}
\newcommand{\ControlFlowTok}[1]{\textcolor[rgb]{0.13,0.29,0.53}{\textbf{#1}}}
\newcommand{\DataTypeTok}[1]{\textcolor[rgb]{0.13,0.29,0.53}{#1}}
\newcommand{\DecValTok}[1]{\textcolor[rgb]{0.00,0.00,0.81}{#1}}
\newcommand{\DocumentationTok}[1]{\textcolor[rgb]{0.56,0.35,0.01}{\textbf{\textit{#1}}}}
\newcommand{\ErrorTok}[1]{\textcolor[rgb]{0.64,0.00,0.00}{\textbf{#1}}}
\newcommand{\ExtensionTok}[1]{#1}
\newcommand{\FloatTok}[1]{\textcolor[rgb]{0.00,0.00,0.81}{#1}}
\newcommand{\FunctionTok}[1]{\textcolor[rgb]{0.13,0.29,0.53}{\textbf{#1}}}
\newcommand{\ImportTok}[1]{#1}
\newcommand{\InformationTok}[1]{\textcolor[rgb]{0.56,0.35,0.01}{\textbf{\textit{#1}}}}
\newcommand{\KeywordTok}[1]{\textcolor[rgb]{0.13,0.29,0.53}{\textbf{#1}}}
\newcommand{\NormalTok}[1]{#1}
\newcommand{\OperatorTok}[1]{\textcolor[rgb]{0.81,0.36,0.00}{\textbf{#1}}}
\newcommand{\OtherTok}[1]{\textcolor[rgb]{0.56,0.35,0.01}{#1}}
\newcommand{\PreprocessorTok}[1]{\textcolor[rgb]{0.56,0.35,0.01}{\textit{#1}}}
\newcommand{\RegionMarkerTok}[1]{#1}
\newcommand{\SpecialCharTok}[1]{\textcolor[rgb]{0.81,0.36,0.00}{\textbf{#1}}}
\newcommand{\SpecialStringTok}[1]{\textcolor[rgb]{0.31,0.60,0.02}{#1}}
\newcommand{\StringTok}[1]{\textcolor[rgb]{0.31,0.60,0.02}{#1}}
\newcommand{\VariableTok}[1]{\textcolor[rgb]{0.00,0.00,0.00}{#1}}
\newcommand{\VerbatimStringTok}[1]{\textcolor[rgb]{0.31,0.60,0.02}{#1}}
\newcommand{\WarningTok}[1]{\textcolor[rgb]{0.56,0.35,0.01}{\textbf{\textit{#1}}}}
\usepackage{graphicx}
\makeatletter
\newsavebox\pandoc@box
\newcommand*\pandocbounded[1]{% scales image to fit in text height/width
  \sbox\pandoc@box{#1}%
  \Gscale@div\@tempa{\textheight}{\dimexpr\ht\pandoc@box+\dp\pandoc@box\relax}%
  \Gscale@div\@tempb{\linewidth}{\wd\pandoc@box}%
  \ifdim\@tempb\p@<\@tempa\p@\let\@tempa\@tempb\fi% select the smaller of both
  \ifdim\@tempa\p@<\p@\scalebox{\@tempa}{\usebox\pandoc@box}%
  \else\usebox{\pandoc@box}%
  \fi%
}
% Set default figure placement to htbp
\def\fps@figure{htbp}
\makeatother
\setlength{\emergencystretch}{3em} % prevent overfull lines
\providecommand{\tightlist}{%
  \setlength{\itemsep}{0pt}\setlength{\parskip}{0pt}}
\usepackage{bookmark}
\IfFileExists{xurl.sty}{\usepackage{xurl}}{} % add URL line breaks if available
\urlstyle{same}
\hypersetup{
  pdftitle={SEATRAC Introduction to R: Session 2},
  pdfauthor={Ling Wei \& Shuyi Ma \textbar{} adapted from the materials by Kim Dill-McFarland's BIGslu tidyverse workshop},
  hidelinks,
  pdfcreator={LaTeX via pandoc}}

\title{SEATRAC Introduction to R: Session 2}
\author{Ling Wei \& Shuyi Ma \textbar{} adapted from the materials by
Kim Dill-McFarland's
\href{https://bigslu.github.io/workshops/2022.08.15_R.tidyverse.workshop/2_tidyverse.html\#2_Introduction_to_the_tidyverse}{BIGslu
tidyverse workshop}}
\date{version January 26, 2026}

\begin{document}
\maketitle

\section{Overview}\label{overview}

In this session, we will introduce you to data manipulation and
visualization using \texttt{tidyverse}, a collection of R packages for
data science. We will cover the following topics:

\begin{itemize}
\tightlist
\item
  retrieving data frame specifications
\item
  chaining sequences of operations (pipes)
\item
  manipulating data frames (select, add, rearrange columns/rows)
\item
  calculating statistics for the data
\item
  data visualization (scatter plot, box plot)
\end{itemize}

If you would like to see an overview of commands we are discussing in
this workshop, this
\href{https://ma-lab-seattle-childrens-cgidr.github.io/SEATRAC_Rintro/Session2_cheatsheet2.html}{cheatsheet}
provides them in a concise view.

\section{Data specifications}\label{data-specifications}

\texttt{tidyverse} is not loaded on R by default. We will load
\texttt{tidyverse} first. If the package has not been installed yet,
then we will need to use the \texttt{install.packages()} function to
install first.

\begin{Shaded}
\begin{Highlighting}[]
\FunctionTok{library}\NormalTok{(tidyverse)}
\end{Highlighting}
\end{Shaded}

\begin{verbatim}
## Warning: package 'readr' was built under R version 4.5.2
\end{verbatim}

\begin{verbatim}
## -- Attaching core tidyverse packages ------------------------ tidyverse 2.0.0 --
## v dplyr     1.1.4     v readr     2.1.6
## v forcats   1.0.1     v stringr   1.6.0
## v ggplot2   3.5.2     v tibble    3.3.0
## v lubridate 1.9.4     v tidyr     1.3.1
## v purrr     1.2.0     
## -- Conflicts ------------------------------------------ tidyverse_conflicts() --
## x dplyr::filter() masks stats::filter()
## x dplyr::lag()    masks stats::lag()
## i Use the conflicted package (<http://conflicted.r-lib.org/>) to force all conflicts to become errors
\end{verbatim}

We will use the toy heart disease dataset to demonstrate the use of
tidyverse for data manipulation and visualization in this session. Here
is the link to download:
\href{https://github.com/Ma-Lab-Seattle-Childrens-CGIDR/SEATRAC_Rintro/blob/main/Heart_Disease_100sampled.csv}{Heart\_Disease\_100sampled.csv}

(Note: this spreadsheet is a sub-sampled dataset from
\href{https://www.kaggle.com/datasets/oktayrdeki/heart-disease/data}{kaggle}.)

To make things a bit easier, we can set our working directory to where
the dataset and the R script are. To do so, navigate to Session
-\textgreater{} Set Working Directory -\textgreater{} Choose Directory,
and then select its location on your computer.

\begin{Shaded}
\begin{Highlighting}[]
\NormalTok{data }\OtherTok{=} \FunctionTok{read\_csv}\NormalTok{(}\StringTok{"Heart\_Disease\_100sampled.csv"}\NormalTok{)}
\end{Highlighting}
\end{Shaded}

\begin{verbatim}
## Rows: 100 Columns: 22
## -- Column specification --------------------------------------------------------
## Delimiter: ","
## chr (13): ptID, Gender, Exercise Habits, Smoking, Family Heart Disease, Diab...
## dbl  (9): Age, Blood Pressure, Cholesterol Level, BMI, Sleep Hours, Triglyce...
## 
## i Use `spec()` to retrieve the full column specification for this data.
## i Specify the column types or set `show_col_types = FALSE` to quiet this message.
\end{verbatim}

Let's first check the dimension of the data.

\begin{Shaded}
\begin{Highlighting}[]
\FunctionTok{dim}\NormalTok{(data)}
\end{Highlighting}
\end{Shaded}

\begin{verbatim}
## [1] 100  22
\end{verbatim}

We can extract and display the full column specification for the data.

\begin{Shaded}
\begin{Highlighting}[]
\FunctionTok{spec}\NormalTok{(data)}
\end{Highlighting}
\end{Shaded}

\begin{verbatim}
## cols(
##   ptID = col_character(),
##   Age = col_double(),
##   Gender = col_character(),
##   `Blood Pressure` = col_double(),
##   `Cholesterol Level` = col_double(),
##   `Exercise Habits` = col_character(),
##   Smoking = col_character(),
##   `Family Heart Disease` = col_character(),
##   Diabetes = col_character(),
##   BMI = col_double(),
##   `High Blood Pressure` = col_character(),
##   `Low HDL Cholesterol` = col_character(),
##   `High LDL Cholesterol` = col_character(),
##   `Alcohol Consumption` = col_character(),
##   `Stress Level` = col_character(),
##   `Sleep Hours` = col_double(),
##   `Sugar Consumption` = col_character(),
##   `Triglyceride Level` = col_double(),
##   `Fasting Blood Sugar` = col_double(),
##   `CRP Level` = col_double(),
##   `Homocysteine Level` = col_double(),
##   `Heart Disease Status` = col_character()
## )
\end{verbatim}

This will tell you the data types (e.g., character, double, integer,
logical) for each column in the data frame.

If you'd like to have a better idea about the structure or content of
the data:

\begin{Shaded}
\begin{Highlighting}[]
\FunctionTok{View}\NormalTok{(data)}
\end{Highlighting}
\end{Shaded}

This will open a new tab to display a preview of the data (maximum 1,000
rows and 50 columns) in a table format, where you can browse, sort, or
search within the data, but you cannot edit the data (read-only) through
this interface.

\section{Columns}\label{columns}

We will work on data frame columns first. To pick specific columns from
a data frame, you can use the \texttt{select()} function in
\texttt{dplyr} package and refer to columns by name, position, or
specific criteria.

To select a single column by name:

\begin{Shaded}
\begin{Highlighting}[]
\NormalTok{data\_ptID }\OtherTok{=} \FunctionTok{select}\NormalTok{(data, ptID)}
\end{Highlighting}
\end{Shaded}

To select multiple columns by name:

\begin{Shaded}
\begin{Highlighting}[]
\NormalTok{data\_multipleCols }\OtherTok{=} \FunctionTok{select}\NormalTok{(data, ptID, BMI, }\StringTok{\textasciigrave{}}\AttributeTok{CRP Level}\StringTok{\textasciigrave{}}\NormalTok{)}
\end{Highlighting}
\end{Shaded}

Column names can be changed as well (e.g., you can remove spaces or use
a more descriptive name). You can use the \texttt{rename()} function in
\texttt{dplyr} package to do it.

\begin{Shaded}
\begin{Highlighting}[]
\NormalTok{data\_renameAge }\OtherTok{=} \FunctionTok{rename}\NormalTok{(data, }\StringTok{\textasciigrave{}}\AttributeTok{Age (Year)}\StringTok{\textasciigrave{}} \OtherTok{=}\NormalTok{ Age)}
\end{Highlighting}
\end{Shaded}

This will add the unit ``Year'' for the ``Age'' column.

You can also create new columns or modify existing columns in a data
frame using the \texttt{mutate()} function in \texttt{dplyr} package.
For instance, if you want to add a new column to convert age in year to
age in month:

\begin{Shaded}
\begin{Highlighting}[]
\NormalTok{data\_addCol }\OtherTok{=} \FunctionTok{mutate}\NormalTok{(data\_renameAge, }\StringTok{\textasciigrave{}}\AttributeTok{Age (Month)}\StringTok{\textasciigrave{}} \OtherTok{=} \StringTok{\textasciigrave{}}\AttributeTok{Age (Year)}\StringTok{\textasciigrave{}} \SpecialCharTok{*} \DecValTok{12}\NormalTok{)}
\end{Highlighting}
\end{Shaded}

This will add a new column named ``Age (Month)'' to the end of the
original data frame.

If you'd like to move this new column right after the ``Age (Year)''
column, you can use the \texttt{relocate()} function in \texttt{dplyr}
package to do it.

\begin{Shaded}
\begin{Highlighting}[]
\NormalTok{data\_addCol\_rearrange }\OtherTok{=} \FunctionTok{relocate}\NormalTok{(data\_addCol, }\StringTok{\textasciigrave{}}\AttributeTok{Age (Month)}\StringTok{\textasciigrave{}}\NormalTok{, }\AttributeTok{.after =} \StringTok{\textasciigrave{}}\AttributeTok{Age (Year)}\StringTok{\textasciigrave{}}\NormalTok{)}
\end{Highlighting}
\end{Shaded}

\section{Rows}\label{rows}

Let's work on subsetting data frame rows using functions from the dplyr
package. Similar to column manipulations, you can select specific rows
from a data frame by position, value of a variable, or specific
criteria.

For example, if you want to extract all the measurements for patients
who smoke, you can use the \texttt{filter()} function to keep rows in
the data frame that match this criteria:

\begin{Shaded}
\begin{Highlighting}[]
\NormalTok{data\_smoke }\OtherTok{=} \FunctionTok{filter}\NormalTok{(data, Smoking }\SpecialCharTok{==} \StringTok{"Yes"}\NormalTok{)}
\end{Highlighting}
\end{Shaded}

You can also add more criteria to the selection using \texttt{\&} (and)
or \texttt{\textbar{}} (or) operators. For instance, if you want to
extract data from patients who not only smoke but also do not exercise
regularly, simply add that condition into the argument like this:

\begin{Shaded}
\begin{Highlighting}[]
\NormalTok{data\_smoke\_noExcercise }\OtherTok{=} \FunctionTok{filter}\NormalTok{(data, Smoking }\SpecialCharTok{==} \StringTok{"Yes"} \SpecialCharTok{\&} \StringTok{\textasciigrave{}}\AttributeTok{Exercise Habits}\StringTok{\textasciigrave{}} \SpecialCharTok{==} \StringTok{"Low"}\NormalTok{)}
\end{Highlighting}
\end{Shaded}

You can also subset rows by specifying their row indexes in a data frame
using the \texttt{slice()} function in \texttt{dplyr} package.

\begin{Shaded}
\begin{Highlighting}[]
\NormalTok{data\_subsetRows }\OtherTok{=} \FunctionTok{slice}\NormalTok{(data, }\DecValTok{20}\SpecialCharTok{:}\DecValTok{50}\NormalTok{)}
\end{Highlighting}
\end{Shaded}

This will give you all the rows between and including row 20 and row 50
in the example dataset.

You can also reorder the rows of the entire data frame based on values
of one or more columns using the \texttt{arrange()} function in
\texttt{dplyr} package. The rows are sorted in ascending order by
default, but you can use the \texttt{desc()} helper function to sort in
descending order.

\begin{Shaded}
\begin{Highlighting}[]
\NormalTok{data\_asc }\OtherTok{=} \FunctionTok{arrange}\NormalTok{(data, Age) }\CommentTok{\# ascending order}
\NormalTok{data\_desc }\OtherTok{=} \FunctionTok{arrange}\NormalTok{(data, }\FunctionTok{desc}\NormalTok{(Age)) }\CommentTok{\# descending order}
\end{Highlighting}
\end{Shaded}

\section{Statistics}\label{statistics}

To calculate some summary statistics for the data, you can use the
\texttt{summarise()} function in \texttt{dplyr} package, which reduces
multiple values from a data frame into a single or multiple summary
statistics.

\begin{Shaded}
\begin{Highlighting}[]
\FunctionTok{summarise}\NormalTok{(data, }\FunctionTok{mean}\NormalTok{(BMI), }\FunctionTok{median}\NormalTok{(BMI)) }\CommentTok{\# mean and median}
\end{Highlighting}
\end{Shaded}

\begin{verbatim}
## # A tibble: 1 x 2
##   `mean(BMI)` `median(BMI)`
##         <dbl>         <dbl>
## 1        29.0          29.5
\end{verbatim}

\begin{Shaded}
\begin{Highlighting}[]
\FunctionTok{summarise}\NormalTok{(data, }\FunctionTok{sd}\NormalTok{(BMI), }\FunctionTok{var}\NormalTok{(BMI)) }\CommentTok{\# standard deviation and variance}
\end{Highlighting}
\end{Shaded}

\begin{verbatim}
## # A tibble: 1 x 2
##   `sd(BMI)` `var(BMI)`
##       <dbl>      <dbl>
## 1      5.97       35.7
\end{verbatim}

\begin{Shaded}
\begin{Highlighting}[]
\FunctionTok{summarise}\NormalTok{(data, }\FunctionTok{min}\NormalTok{(BMI), }\FunctionTok{max}\NormalTok{(BMI)) }\CommentTok{\# minimum and maximum}
\end{Highlighting}
\end{Shaded}

\begin{verbatim}
## # A tibble: 1 x 2
##   `min(BMI)` `max(BMI)`
##        <dbl>      <dbl>
## 1       18.2       39.8
\end{verbatim}

\subsection{Getting help}\label{getting-help}

If you want to learn more about the \texttt{dplyr} package, please refer
to this ``Introduction to dplyr'' guide, which provides a comprehensive
and detailed description of the package (including examples).

\begin{Shaded}
\begin{Highlighting}[]
\FunctionTok{vignette}\NormalTok{(}\StringTok{"dplyr"}\NormalTok{)}
\end{Highlighting}
\end{Shaded}

\begin{verbatim}
## starting httpd help server ... done
\end{verbatim}

\subsection{Bonus}\label{bonus}

To select all columns between two specific columns in a data frame, you
can specify the names of the first and the last columns that you'd like
to include.

\begin{Shaded}
\begin{Highlighting}[]
\NormalTok{data\_colsRange }\OtherTok{=} \FunctionTok{select}\NormalTok{(data, ptID}\SpecialCharTok{:}\NormalTok{BMI)}
\end{Highlighting}
\end{Shaded}

This will give you all the columns in the example dataset between (and
including) columns ``ptID'' and ``BMI''.

You can also use the \texttt{select()} function to exclude one
particular column in a data frame. For instance, if you want to retrieve
all the measurements in the example dataset except for sleep hours, you
can simply specify its column name with the \texttt{!}:

\begin{Shaded}
\begin{Highlighting}[]
\NormalTok{data\_excludeCol }\OtherTok{=} \FunctionTok{select}\NormalTok{(data, }\SpecialCharTok{!}\StringTok{\textasciigrave{}}\AttributeTok{Sleep Hours}\StringTok{\textasciigrave{}}\NormalTok{)}
\end{Highlighting}
\end{Shaded}

If you want to get the first or the last number of rows from a data
frame, use the \texttt{slice\_head()} or \texttt{slice\_tail()}
function, respectively, and specify the number of rows you'd like to
subset.

\begin{Shaded}
\begin{Highlighting}[]
\NormalTok{data\_headRows }\OtherTok{=} \FunctionTok{slice\_head}\NormalTok{(data, }\AttributeTok{n=}\DecValTok{5}\NormalTok{)}
\NormalTok{data\_tailRows }\OtherTok{=} \FunctionTok{slice\_tail}\NormalTok{(data, }\AttributeTok{n=}\DecValTok{5}\NormalTok{)}
\end{Highlighting}
\end{Shaded}

You can also randomly select rows from a data frame. This can be done
with the \texttt{slice\_sample()} function in \texttt{dplyr} package.
For instance, if you want to get data from 10 patients that are randomly
sampled from the 100 patients in the example dataset:

\begin{Shaded}
\begin{Highlighting}[]
\NormalTok{data\_randomRows }\OtherTok{=} \FunctionTok{slice\_sample}\NormalTok{(data, }\AttributeTok{n =} \DecValTok{10}\NormalTok{)}
\end{Highlighting}
\end{Shaded}

Row selections can also be done based on the values of a particular
variable in a data frame. For example, you can use the
\texttt{slice\_min()} function to extract data from the five patients
with the lowest BMI:

\begin{Shaded}
\begin{Highlighting}[]
\NormalTok{data\_lowBMI }\OtherTok{=} \FunctionTok{slice\_min}\NormalTok{(data, }\AttributeTok{order\_by =}\NormalTok{ BMI, }\AttributeTok{n=}\DecValTok{5}\NormalTok{)}
\end{Highlighting}
\end{Shaded}

Or if you want to get data from the five oldest patients, use the
\texttt{slice\_max()} function:

\begin{Shaded}
\begin{Highlighting}[]
\NormalTok{data\_oldPt }\OtherTok{=} \FunctionTok{slice\_max}\NormalTok{(data, }\AttributeTok{order\_by =}\NormalTok{ Age, }\AttributeTok{n=}\DecValTok{5}\NormalTok{)}
\end{Highlighting}
\end{Shaded}

\section{Pipes}\label{pipes}

Oftentimes your code contains multiple nested parenthesis, which makes
it hard to read and understand. The pipe operator
\texttt{\%\textgreater{}\%} allows you to transform complex nested
operations into a linear sequence of operations, significantly increase
the readability of the code.

Another potential advantage is that if you use a pipe to perform a
series of operations, then you also won't need to store intermediate
outputs as separate variables. This might make it easier to keep track
of which output variables you actually care about.

For example, to select a single column from the data frame using a pipe:

\begin{Shaded}
\begin{Highlighting}[]
\NormalTok{data }\SpecialCharTok{\%\textgreater{}\%}
  \FunctionTok{select}\NormalTok{(ptID)}
\end{Highlighting}
\end{Shaded}

\begin{verbatim}
## # A tibble: 100 x 1
##    ptID   
##    <chr>  
##  1 pt 7186
##  2 pt 7572
##  3 pt 2570
##  4 pt 8253
##  5 pt 5216
##  6 pt 3775
##  7 pt 7529
##  8 pt 8049
##  9 pt 1862
## 10 pt 8058
## # i 90 more rows
\end{verbatim}

We will write the code in piped format in the rest of this workshop.

\section{Data visualization}\label{data-visualization}

We will make some simple plots from the toy dataset using
\texttt{ggplot2}, a widely used R package for data visualization.
\texttt{ggplot2} allows you to build complex plots layer by layer
through the \texttt{+} operator. You will first use \texttt{ggplot()} to
define the plot and aesthetic mappings (e.g., axes, color, size). You
will then add layers to define the geometry (e.g., points, lines, bars)
of the plot. You can also add layers to create subplots, add title and
axis labels, and save the plot.

\subsection{scatter plot}\label{scatter-plot}

We will start with scatter plots. We will make a basic scatter plot of
``Fasting Blood Sugar'' against ``CRP Level'' using the default
settings.

\begin{Shaded}
\begin{Highlighting}[]
\NormalTok{data }\SpecialCharTok{\%\textgreater{}\%}
  \FunctionTok{ggplot}\NormalTok{(}\FunctionTok{aes}\NormalTok{(}\AttributeTok{x =} \StringTok{\textasciigrave{}}\AttributeTok{Fasting Blood Sugar}\StringTok{\textasciigrave{}}\NormalTok{, }\AttributeTok{y =} \StringTok{\textasciigrave{}}\AttributeTok{CRP Level}\StringTok{\textasciigrave{}}\NormalTok{)) }\SpecialCharTok{+}
  \FunctionTok{geom\_point}\NormalTok{()}
\end{Highlighting}
\end{Shaded}

\pandocbounded{\includegraphics[keepaspectratio]{Session2_tidyR_Plotting.v1_files/figure-latex/basic scatter plot-1.pdf}}

You can color code the data points by gender by adding a \texttt{color}
argument. You can also modify the appearance of the plot. For instance,
we can change it to a white background with gray grid lines using the
\texttt{theme\_bw()} function. To adjust axis labels or to add a title
to the plot, you can use the \texttt{labs()} and \texttt{theme()}
functions.

\begin{Shaded}
\begin{Highlighting}[]
\NormalTok{data }\SpecialCharTok{\%\textgreater{}\%}
  \FunctionTok{ggplot}\NormalTok{(}\FunctionTok{aes}\NormalTok{(}\AttributeTok{x =} \StringTok{\textasciigrave{}}\AttributeTok{Fasting Blood Sugar}\StringTok{\textasciigrave{}}\NormalTok{, }\AttributeTok{y =} \StringTok{\textasciigrave{}}\AttributeTok{CRP Level}\StringTok{\textasciigrave{}}\NormalTok{, }\AttributeTok{color =}\NormalTok{ Gender)) }\SpecialCharTok{+}
  \FunctionTok{geom\_point}\NormalTok{() }\SpecialCharTok{+}
  \FunctionTok{theme\_bw}\NormalTok{() }\SpecialCharTok{+}
  \FunctionTok{labs}\NormalTok{(}\AttributeTok{x =} \StringTok{"Fasting Blood Sugar [mg/dL]"}\NormalTok{, }\AttributeTok{y =} \StringTok{"CRP Level [mg/L]"}\NormalTok{, }\AttributeTok{title =} \StringTok{"Fasting blood sugar VS CRP level"}\NormalTok{) }\SpecialCharTok{+}
  \FunctionTok{theme}\NormalTok{(}
    \AttributeTok{axis.title.x =} \FunctionTok{element\_text}\NormalTok{(}\AttributeTok{size =} \DecValTok{14}\NormalTok{, }\AttributeTok{color =} \StringTok{"red"}\NormalTok{),}
    \AttributeTok{axis.title.y =} \FunctionTok{element\_text}\NormalTok{(}\AttributeTok{size =} \DecValTok{14}\NormalTok{, }\AttributeTok{color =} \StringTok{"blue"}\NormalTok{),}
    \AttributeTok{plot.title =} \FunctionTok{element\_text}\NormalTok{(}\AttributeTok{hjust =} \FloatTok{0.5}\NormalTok{, }\AttributeTok{face =} \StringTok{"bold"}\NormalTok{, }\AttributeTok{size =} \DecValTok{18}\NormalTok{, }\AttributeTok{color =} \StringTok{"black"}\NormalTok{)}
\NormalTok{    )}
\end{Highlighting}
\end{Shaded}

\pandocbounded{\includegraphics[keepaspectratio]{Session2_tidyR_Plotting.v1_files/figure-latex/scatter plot settings-1.pdf}}

To help visualize patterns and trends in the data, you can use the
\texttt{geom\_smooth()} function to fit a model to the data and overlay
a trend line to the plot. It also displays shaded 95\% confidence
interval by default (you can hide or adjust it in the arguments). You
can specify the function to use for fitting the model (smoothing
function) via the \texttt{method} argument.

For example, if you want to add linear regression lines to the plot
above, you can add a layer of smoothing function \texttt{geom\_smooth()}
and put \texttt{lm} in the ``method'' argument.

\begin{Shaded}
\begin{Highlighting}[]
\NormalTok{data }\SpecialCharTok{\%\textgreater{}\%}
  \FunctionTok{ggplot}\NormalTok{(}\FunctionTok{aes}\NormalTok{(}\AttributeTok{x =} \StringTok{\textasciigrave{}}\AttributeTok{Fasting Blood Sugar}\StringTok{\textasciigrave{}}\NormalTok{, }\AttributeTok{y =} \StringTok{\textasciigrave{}}\AttributeTok{Homocysteine Level}\StringTok{\textasciigrave{}}\NormalTok{, }\AttributeTok{color =}\NormalTok{ Gender)) }\SpecialCharTok{+}
  \FunctionTok{geom\_point}\NormalTok{() }\SpecialCharTok{+}
  \FunctionTok{theme\_bw}\NormalTok{() }\SpecialCharTok{+}
  \FunctionTok{labs}\NormalTok{(}\AttributeTok{x =} \StringTok{"Fasting Blood Sugar [mg/dL]"}\NormalTok{, }\AttributeTok{y =} \StringTok{"CRP Level [mg/L]"}\NormalTok{, }\AttributeTok{title =} \StringTok{"Fasting blood sugar VS CRP level"}\NormalTok{) }\SpecialCharTok{+}
  \FunctionTok{theme}\NormalTok{(}
    \AttributeTok{axis.title.x =} \FunctionTok{element\_text}\NormalTok{(}\AttributeTok{size =} \DecValTok{14}\NormalTok{, }\AttributeTok{color =} \StringTok{"red"}\NormalTok{),}
    \AttributeTok{axis.title.y =} \FunctionTok{element\_text}\NormalTok{(}\AttributeTok{size =} \DecValTok{14}\NormalTok{, }\AttributeTok{color =} \StringTok{"blue"}\NormalTok{),}
    \AttributeTok{plot.title =} \FunctionTok{element\_text}\NormalTok{(}\AttributeTok{hjust =} \FloatTok{0.5}\NormalTok{, }\AttributeTok{face =} \StringTok{"bold"}\NormalTok{, }\AttributeTok{size =} \DecValTok{18}\NormalTok{, }\AttributeTok{color =} \StringTok{"black"}\NormalTok{)}
\NormalTok{    ) }\SpecialCharTok{+}
  \FunctionTok{geom\_smooth}\NormalTok{(}\AttributeTok{method =} \StringTok{"lm"}\NormalTok{, }\AttributeTok{se =} \ConstantTok{TRUE}\NormalTok{)}
\end{Highlighting}
\end{Shaded}

\begin{verbatim}
## `geom_smooth()` using formula = 'y ~ x'
\end{verbatim}

\pandocbounded{\includegraphics[keepaspectratio]{Session2_tidyR_Plotting.v1_files/figure-latex/regression line-1.pdf}}

Since we group the data by gender here, the fitting is applied to female
and male patient data separately.

\subsection{box-and-whisker plot}\label{box-and-whisker-plot}

We will make some box-and-whisker plots to visualize the distribution of
data between groups. We will start with a basic box plot on triglyceride
level for patients with or without heart disease. We will use
\texttt{geom\_boxplot()} to define the box plot. If you want the y axis
to show the full data range and hide the outliers, specify in the
argument with \texttt{outlier.shape\ =\ NA}. You can use the
\texttt{geom\_jitter()} function to display individual data points,
which also adds a small amount of random noise to the position
(horizontal in this case) of each point for easy visualization.

\begin{Shaded}
\begin{Highlighting}[]
\NormalTok{data }\SpecialCharTok{\%\textgreater{}\%}
  \FunctionTok{ggplot}\NormalTok{(}\FunctionTok{aes}\NormalTok{(}\AttributeTok{x =} \StringTok{\textasciigrave{}}\AttributeTok{Heart Disease Status}\StringTok{\textasciigrave{}}\NormalTok{, }\AttributeTok{y =} \StringTok{\textasciigrave{}}\AttributeTok{Triglyceride Level}\StringTok{\textasciigrave{}}\NormalTok{)) }\SpecialCharTok{+}
  \FunctionTok{geom\_boxplot}\NormalTok{(}\AttributeTok{outlier.shape =} \ConstantTok{NA}\NormalTok{) }\SpecialCharTok{+}
  \FunctionTok{geom\_jitter}\NormalTok{() }\SpecialCharTok{+}
  \FunctionTok{theme\_bw}\NormalTok{()}
\end{Highlighting}
\end{Shaded}

\pandocbounded{\includegraphics[keepaspectratio]{Session2_tidyR_Plotting.v1_files/figure-latex/basic box plot-1.pdf}}

By default, \texttt{geom\_boxplot()} does not include error bars (or
whisker end-bars) in the plot. You can call the \texttt{stat\_boxplot()}
function to add horizontal lines at the end of whiskers.

\begin{Shaded}
\begin{Highlighting}[]
\NormalTok{data }\SpecialCharTok{\%\textgreater{}\%}
  \FunctionTok{ggplot}\NormalTok{(}\FunctionTok{aes}\NormalTok{(}\AttributeTok{x =} \StringTok{\textasciigrave{}}\AttributeTok{Heart Disease Status}\StringTok{\textasciigrave{}}\NormalTok{, }\AttributeTok{y =} \StringTok{\textasciigrave{}}\AttributeTok{Triglyceride Level}\StringTok{\textasciigrave{}}\NormalTok{)) }\SpecialCharTok{+}
  \FunctionTok{geom\_boxplot}\NormalTok{(}\AttributeTok{outlier.shape =} \ConstantTok{NA}\NormalTok{) }\SpecialCharTok{+}
  \FunctionTok{stat\_boxplot}\NormalTok{(}\AttributeTok{geom =} \StringTok{"errorbar"}\NormalTok{, }\AttributeTok{width =} \FloatTok{0.2}\NormalTok{) }\SpecialCharTok{+}
  \FunctionTok{geom\_jitter}\NormalTok{() }\SpecialCharTok{+}
  \FunctionTok{theme\_bw}\NormalTok{()}
\end{Highlighting}
\end{Shaded}

\pandocbounded{\includegraphics[keepaspectratio]{Session2_tidyR_Plotting.v1_files/figure-latex/error bars-1.pdf}}
You can also color code the boxes, bars, and individual data points by
groups. To do so, simply add an argument \texttt{color} in the
aesthetics function \texttt{aes()}.

\begin{Shaded}
\begin{Highlighting}[]
\NormalTok{data }\SpecialCharTok{\%\textgreater{}\%}
  \FunctionTok{ggplot}\NormalTok{(}\FunctionTok{aes}\NormalTok{(}\AttributeTok{x =} \StringTok{\textasciigrave{}}\AttributeTok{Heart Disease Status}\StringTok{\textasciigrave{}}\NormalTok{, }\AttributeTok{y =} \StringTok{\textasciigrave{}}\AttributeTok{Triglyceride Level}\StringTok{\textasciigrave{}}\NormalTok{, }\AttributeTok{color =} \StringTok{\textasciigrave{}}\AttributeTok{Heart Disease Status}\StringTok{\textasciigrave{}}\NormalTok{)) }\SpecialCharTok{+}
  \FunctionTok{geom\_boxplot}\NormalTok{(}\AttributeTok{outlier.shape =} \ConstantTok{NA}\NormalTok{) }\SpecialCharTok{+}
  \FunctionTok{stat\_boxplot}\NormalTok{(}\AttributeTok{geom =} \StringTok{"errorbar"}\NormalTok{, }\AttributeTok{width =} \FloatTok{0.2}\NormalTok{) }\SpecialCharTok{+}
  \FunctionTok{geom\_jitter}\NormalTok{() }\SpecialCharTok{+}
  \FunctionTok{theme\_bw}\NormalTok{()}
\end{Highlighting}
\end{Shaded}

\pandocbounded{\includegraphics[keepaspectratio]{Session2_tidyR_Plotting.v1_files/figure-latex/colored box plot-1.pdf}}

\end{document}
